\documentclass[11pt]{article}
\usepackage{amsmath,amssymb,amsthm}
\newtheorem{theorem}{Theorem}
\title{Exact progress from a clearance band with different actuator limits}
\begin{document}
\maketitle

Requiring the two robots to move in lockstep makes them inherit the
slower robot's speed and the slower robot's acceleration. Allowing a
small, physically bounded change of gap can recover some of that
progress, but does not make the individually feasible targets joint.

\paragraph{Law.}
Robot 1 has acceleration bound \(a>0\) and speed bound \(W\).
Robot 2 has acceleration bound \(A\ge a\) and speed bound
\(V>0\), with \(W\ge V\). Velocity is absolutely continuous.
Subtract the initial positions and let \(x_i(t)\) be their advances.
Both start at \(x_i=0\) at rest, and both must be at the same advance
\(D\) and at rest at time \(T\). The allowable gap variation is
\(|x_2-x_1|\le\eta\) at every time. For an initial physical gap \(g\)
this enforces a lower clearance \(g-\eta>0\) and an upper link
distance \(g+\eta\). The test has no obstacles, staggered notice,
unknown coefficients, or radio guarantee beyond geometry.
The same kinematic limits apply to any independent comparison.

\begin{theorem}
Let \(B(t;a,v)=\int_0^t\min(a u,a(t-u),v)\,du\), the standalone
rest-to-rest travel bound, and write \(\tau=V/a\). If \(T\le2\tau\),
the exact maximum is \(aT^2/4\). Otherwise the exact maximum is
\[
 \min\bigl(B(T;a,W),\ B(T;A,V),\ VT-V^2/a+2\eta\bigr).
\]
\end{theorem}
\begin{proof}
For \(T\le2\tau\), robot 1's triangular speed profile has peak
at most \(V\) and both can realize it with the same motion. Its
own acceleration bound prevents any larger advance.

Now let \(H=T-2\tau>0\) and \(E=VT-V^2/a\). At time \(\tau\),
acceleration from rest gives \(x_1(\tau)\le a\tau^2/2\), so
\(x_2(\tau)\le a\tau^2/2+\eta\). At time \(T-\tau\), the
terminal rest gives \(x_1(T-\tau)\ge D-a\tau^2/2\), so
\(x_2(T-\tau)\ge D-a\tau^2/2-\eta\). The middle displacement of
robot 2 is at most \(VH\). Combining these yields \(D\le E+2\eta\).
The argument covers every admissible velocity, including reversal.
Each robot also has its own rest-to-rest upper bound.

Let \(e\) be half the difference between the claimed maximum
and \(E\); \(0\le e\le\eta\). The other two terms in the
minimum ensure
\[
        2e\le V\tau-V^2/A,\qquad
        2e\le B(H;a,W-V).
\]
Indeed robot 2's standalone optimum is \(VT-V^2/A\),
and robot 1's optimum equals \(E+B(H;a,W-V)\):
its speed envelope starts and ends with ramps between zero and
\(V\), and in the middle is \(V\) plus the remaining capped
envelope. The two inequalities describe how much lead robot 2
can develop before the middle, and how much lead robot 1 can
absorb during the middle.

For sharpness, robot 1 accelerates at \(a\) until \(\tau\).
Robot 2 reaches \(V\) at \(s=\tau-2e/V\) with constant
acceleration \(V/s\), then holds \(V\). The first inequality
implies \(s\ge V/A>0\). Robot 2's lead grows monotonically to
\(e\). In the middle, robot 2 travels at \(V\); robot 1 travels
at \(V\) plus the defining profile for \(B(H;a,W-V)\) scaled
to have area \(2e\). Use a zero profile for zero area. The
second inequality lets the scale be at most one. The lead
decreases monotonically to \(-e\). Time-reverse the two initial
speed profiles in the final interval. Both stop together at
\(E+2e\), with the lead approaching zero from below. The speed,
acceleration, and gap bounds hold continuously. Every smaller
nonnegative advance is reached by scaling the paths and velocities.
\end{proof}

\paragraph{Matched tests.}
Take \(a=V=1,A=4,W=3,T=4\). Lockstep permits distance \(3\).
With half-width \(\eta=1/10\), the exact joint maximum rises to
\(16/5\). Both robots can reach \(7/2\) separately under the same
bounds, but doing so together violates the deadline or the gap.
The required half-width for that target is \(1/4\) in this regime.
The replay constructs nine rational velocity profiles and checks the
gap extrema inside every interval, not only at sample points.
An independent time-grid LP with piecewise-linear velocities
agrees on all nine. The proof, rather than the LP's restriction to
that basis, proves the optimum for all continuous paths.

The earlier difference-constraint solution used proportional
velocity and acceleration bounds to move both coordinates along a
single segment. This result instead permits their different velocity
profiles and quantifies the needed gap tolerance. When a standalone limit binds, the construction uses only the gap
slack needed to reach that limit. An independent single-robot distance
is not by itself a joint construction. Standard simultaneous
barrier methods are relevant to viability but are not compared
in operation here. This is a scoped R6 result, not a field
or independent mission test.
\end{document}
